% Generated from common/cover-notes.json by tools/series_cover.py.
\backmatter
\cleardoublepage
\chapter*{About the Cover}
\addcontentsline{toc}{chapter}{About the Cover}
\noindent The integers from 1 to the square of 71 are written in a square spiral, and the primes are the dots. Stanislaw Ulam drew the first such spiral in 1963, doodling during a dull talk, and saw the primes gather along diagonals. Martin Gardner published the figure in Scientific American in 1964. Nothing on the cover explains those diagonals. That is why the spiral stands on this volume, facing the chapter Primes, where the primes stop being a list and become the subject. The drawing was computed by the book's own TikZ cover generator.
